% ============================================================
% Arbol de problemas - Uso problematico de internet (UPI)
% Proyecto de Innovacion 1 - Fase 1 CRISP-ML(Q)
% Compilar en Overleaf con pdfLaTeX
% ============================================================
\documentclass[10pt,letterpaper]{article}
 
\usepackage[utf8]{inputenc}
\usepackage[T1]{fontenc}
\usepackage[spanish,es-noshorthands]{babel}
\usepackage{lmodern}
\usepackage[margin=2.2cm]{geometry}
\usepackage{pdflscape}
\usepackage{tikz}
\usetikzlibrary{positioning,arrows.meta,calc,fit,backgrounds}
\usepackage{xcolor}
\usepackage{booktabs}
\usepackage{tabularx}
\usepackage{enumitem}
\usepackage{array}
\usepackage[hidelinks]{hyperref}
 
% ---------- Colores ----------
\definecolor{efecto}{HTML}{B4462F}
\definecolor{efectosoft}{HTML}{F8E7E2}
\definecolor{causa}{HTML}{1F6F5C}
\definecolor{causasoft}{HTML}{E2F1EC}
\definecolor{tronco}{HTML}{1C2F3B}
\definecolor{ambar}{HTML}{FBF0D9}
 
\setlist[itemize]{leftmargin=1.2em,itemsep=1pt,topsep=2pt,parsep=0pt}
 
% ---------- Estilos TikZ ----------
\tikzset{
  base/.style={rectangle, rounded corners=3pt, align=left, font=\footnotesize,
               inner sep=5pt, text width=5.5cm, minimum height=2.3cm, anchor=north},
  efecto/.style={base, draw=efecto, line width=0.9pt, fill=white},
  efsec/.style={base, draw=efecto!60, fill=efectosoft, minimum height=1.5cm, text width=7.3cm},
  causa/.style={base, draw=causa, line width=0.9pt, fill=white, minimum height=3.2cm},
  problema/.style={rectangle, rounded corners=4pt, fill=tronco, text=white,
               align=center, font=\small\bfseries, text width=23.4cm, inner sep=9pt},
  flecha/.style={-{Stealth[length=2.2mm]}, line width=0.7pt, draw=gray!70},
  titulo/.style={font=\footnotesize\bfseries}
}
 
\title{\textbf{Proyecto de Innovacion 1 \\ \small Fase 1 CRISP-ML(Q): comprension del negocio y los datos}\\[2pt]
\large Uso problematico de internet en ninos y adolescentes}
\author{Rafael Santiago Bedoya, Jose Ricardo Florez \\ Wilson Panesso, Angie Giselle Perez}
\date{}
 
\begin{document}
\maketitle
 
\section*{Contexto}
El reto propuesto por Child Mind Institute y Healthy Brain Network (HBN) busca predecir el
\emph{Severity Impairment Index} (sii), medida estandar del uso problematico de internet (UPI),
a partir de datos de actividad y aptitud fisica de jovenes entre 5 y 22 anos. El sii se deriva
del cuestionario PCIAT: 0 = Ninguno, 1 = Leve, 2 = Moderado, 3 = Severo.
 
% ============================================================
% ARBOL (pagina horizontal)
% ============================================================
\begin{landscape}
\thispagestyle{empty}
\begin{center}
{\large\bfseries Arbol de problemas}\\[4pt]
\begin{tikzpicture}[node distance=0.45cm and 0.35cm]
 
% ----- Problema central -----
\node[problema] (P) {PROBLEMA CENTRAL\\[3pt]
\normalfont\small Deteccion tardia del uso problematico de internet (UPI) en ninos y adolescentes,
porque su medicion depende de evaluaciones clinicas complejas y poco accesibles, y no se
aprovechan indicadores fisicos de bajo costo como senal temprana.};
 
% ----- Efectos directos -----
\node[efecto, above=1.0cm of P.north west, anchor=south west, xshift=0.1cm] (E1)
  {\textcolor{efecto}{\textbf{E1. Intervencion tardia}}
   \begin{itemize}
     \item Habitos digitales nocivos se consolidan antes de detectarse
     \item Se actua cuando ya hay sintomas clinicos
   \end{itemize}};
\node[efecto, right=of E1.north east, anchor=north west] (E2)
  {\textcolor{efecto}{\textbf{E2. Deterioro de la salud mental}}
   \begin{itemize}
     \item Mayor prevalencia de depresion y ansiedad
     \item Alteraciones del sueno
   \end{itemize}};
\node[efecto, right=of E2.north east, anchor=north west] (E3)
  {\textcolor{efecto}{\textbf{E3. Deterioro de la salud fisica}}
   \begin{itemize}
     \item Sedentarismo, mala postura, dieta irregular
     \item Menor aptitud cardiovascular
   \end{itemize}};
\node[efecto, right=of E3.north east, anchor=north west] (E4)
  {\textcolor{efecto}{\textbf{E4. Casos subdiagnosticados}}
   \begin{itemize}
     \item UPI invisible detras de otros diagnosticos
     \item Sin linea base de prevalencia real
   \end{itemize}};
 
% ----- Efectos secundarios -----
\node[efsec, above=0.9cm of E1.north west, anchor=south west, xshift=0cm] (S1)
  {\textbf{Mayores costos} para familias y sistema de salud; politicas publicas sin datos para focalizar};
\node[efsec, right=0.4cm of S1.north east, anchor=north west] (S2)
  {\textbf{Bajo desempeno escolar y social}: menor funcionamiento global (CGAS), aislamiento};
\node[efsec, right=0.4cm of S2.north east, anchor=north west] (S3)
  {\textbf{Mas inequidad}: familias sin acceso a especialistas quedan sin diagnostico};
 
% ----- Causas -----
\node[causa, below=1.0cm of P.south west, anchor=north west, xshift=0.1cm] (C1)
  {\textcolor{causa}{\textbf{C1. Dependencia de evaluaciones clinicas especializadas}}
   \begin{itemize}
     \item PCIAT (20 items) requiere aplicacion profesional
     \item Escasez de especialistas en salud mental infantil
     \item Barreras de acceso, culturales y linguisticas
   \end{itemize}};
\node[causa, right=of C1.north east, anchor=north west] (C2)
  {\textcolor{causa}{\textbf{C2. El UPI no se mide directamente}}
   \begin{itemize}
     \item Se infiere por sus consecuencias (depresion, ansiedad)
     \item Sin tamizaje en colegios ni atencion primaria
     \item Baja conciencia de padres sobre senales de alerta
   \end{itemize}};
\node[causa, right=of C2.north east, anchor=north west] (C3)
  {\textcolor{causa}{\textbf{C3. Subutilizacion de indicadores fisicos accesibles}}
   \begin{itemize}
     \item FitnessGram, bioimpedancia y actigrafia no se conectan con la salud digital
     \item No hay modelos que relacionen actividad, sueno y composicion corporal con el UPI
   \end{itemize}};
\node[causa, right=of C3.north east, anchor=north west] (C4)
  {\textcolor{causa}{\textbf{C4. Datos incompletos y heterogeneos}}
   \begin{itemize}
     \item Muchas mediciones faltantes
     \item sii ausente en parte del entrenamiento
     \item Actigrafia solo en un subconjunto (parquet vs csv)
     \item Clases desbalanceadas
   \end{itemize}};
 
% ----- Flechas: causas -> problema -> efectos -> secundarios -----
\foreach \c in {C1,C2,C3,C4} \draw[flecha, draw=causa] (\c.north) -- (\c.north |- P.south);
\foreach \e in {E1,E2,E3,E4} \draw[flecha, draw=efecto] (\e.south |- P.north) -- (\e.south);
\draw[flecha, draw=efecto!70] (E1.north) -- (S1.south);
\draw[flecha, draw=efecto!70] (E2.north) -- (S1.south);
\draw[flecha, draw=efecto!70] (E4.north) -- (S3.south);
\draw[flecha, draw=efecto!70] (E2.north) -- (S2.south);
\draw[flecha, draw=efecto!70] (E3.north) -- (S2.south);
 
% ----- Etiquetas laterales -----
\node[titulo, text=efecto, rotate=90, left=0.15cm of S1.west, anchor=south] {EFECTOS};
\node[titulo, text=causa, rotate=90, left=0.15cm of C1.west, anchor=south] {CAUSAS};
 
\end{tikzpicture}
\end{center}
\end{landscape}
 
% ============================================================
\section{Planteamiento del problema}
\begin{center}
\colorbox{ambar}{\parbox{0.94\linewidth}{\vspace{6pt}
En ninos y adolescentes de 5 a 22 anos, el uso problematico de internet se detecta tarde porque
su medicion (PCIAT / sii) exige evaluacion clinica especializada, lo que genera barreras de
acceso para muchas familias. Las medidas de aptitud fisica y actividad (FitnessGram,
bioimpedancia, actigrafia de muneca, cuestionarios de actividad y sueno) son baratas y faciles
de tomar, pero no se usan como senal temprana. La consecuencia es una intervencion tardia, con
mayor deterioro de salud mental y fisica e inequidad en el diagnostico.\vspace{6pt}}}
\end{center}
 
\section{Datos que respaldan el problema}
\begin{center}
\small
\begin{tabularx}{\linewidth}{>{\raggedright\arraybackslash}p{4.2cm} >{\raggedright\arraybackslash}p{3.8cm} X}
\toprule
\textbf{Dato} & \textbf{Valor} & \textbf{Por que importa} \\
\midrule
Muestra HBN & $\sim$5.000 jovenes, 5--22 anos & Base clinica con evaluaciones fisicas y de uso de internet \\
Conjunto de prueba oculto & $\sim$3.800 instancias & Escala real sobre la que se evalua el modelo \\
Instrumentos tabulares & 11 & Fuentes heterogeneas que hay que integrar (C3, C4) \\
Actigrafia & Hasta 30 dias, ventanas de 5 s & Senal objetiva de actividad y sueno, solo en un subconjunto \\
Train con sii etiquetado$^{*}$ & $\sim$2.700 de $\sim$3.960 & Justifica aprendizaje semi o no supervisado \\
Distribucion sii$^{*}$ & 0: 58\%, 1: 27\%, 2: 14\%, 3: 1\% & Desbalance fuerte: los casos severos son los mas escasos \\
\bottomrule
\end{tabularx}
\end{center}
{\footnotesize $^{*}$ Cifras aproximadas de la version publica de la competencia en Kaggle; verificar con \texttt{train.csv}.}
 
\section{Puntos de intervencion para IA}
\begin{itemize}
  \item \textbf{C3:} modelo predictivo del sii a partir de datos fisicos, de actividad y sueno (objetivo central del reto).
  \item \textbf{C4:} imputacion de faltantes, ingenieria de variables desde la actigrafia y manejo del desbalance de clases.
  \item \textbf{C1 y C2:} el modelo funciona como tamizaje de bajo costo que prioriza quien necesita evaluacion clinica; no reemplaza el diagnostico.
\end{itemize}
 
\subsection*{Conexion con CRISP-ML(Q)}
\begin{center}
\small
\begin{tabularx}{\linewidth}{>{\raggedright\arraybackslash}p{4.6cm} X}
\toprule
\textbf{Fase} & \textbf{Aplicacion en este proyecto} \\
\midrule
Comprension del negocio y datos & Este arbol, criterios de exito y verificacion de calidad de datos \\
Ingenieria de datos & Integrar csv + parquet, imputar faltantes, balancear clases (C4) \\
Ingenieria y evaluacion del modelo & Metrica ordinal (kappa cuadratico ponderado), robustez y explicabilidad (C3) \\
Garantia de calidad (riesgos) & Sesgo de muestra unica (Nueva York), fuga de informacion del PCIAT, sobreajuste \\
\bottomrule
\end{tabularx}
\end{center}
 
\section{Reflexion}
El arbol muestra que el problema no es la falta de un tratamiento, sino la falta de una senal
temprana y accesible. Las causas C1 y C2 son estructurales y no se resuelven con IA, pero C3 y C4
si: hay datos fisicos baratos que hoy no se aprovechan y un reto tecnico claro en su calidad.
Por eso el proyecto se enfoca en un modelo de tamizaje que prediga el sii con indicadores
fisicos, como apoyo, no sustituto, de la evaluacion clinica.
 
\end{document}
 